\documentclass[sigconf]{acmart}
% Cleaning
\pagestyle{fancy} % removes running headers
\settopmatter{printacmref=false} % Removes citation information below abstract
\renewcommand\footnotetextcopyrightpermission[1]{} % removes footnote with conference information in first column
\pagestyle{plain} % removes running headers
\fancyhead{}
\settopmatter{printacmref=false, printfolios=false}
% Copyright
\setcopyright{none}
\acmDOI{}

\usepackage{listings}
\definecolor{codegray}{rgb}{0.97,0.97,0.97}

\setlength{\parskip}{0pt}
\setlength{\parsep}{0pt}
\setlength{\headsep}{0pt}
\setlength{\topskip}{0pt}
\setlength{\topmargin}{0pt}
\setlength{\topsep}{0pt}
\setlength{\partopsep}{0pt}
%\linespread{0.95}
\usepackage{mdwlist}
\usepackage{amsmath}

\begin{document}
\title{TRIAD: Trusted RDF Interoperable Agent Data}

\author{Asish A. George}
\affiliation{
  \institution{University of Victoria}
  \city{Victoria}
  \state{BC}
  \country{Canada}
}
\email{asishgeorge@uvic.ca}

\author{Mohamed Khaled Abuela}
\affiliation{
  \institution{University of Victoria}
  \city{Victoria}
  \state{BC}
  \country{Canada}
}
\email{mabuela@uvic.ca}

% ABSTRACT: Khaled's text, AI-edited for clarity (Level C). Original in appendix.
\begin{abstract}
Patients can wait and deteriorate while information about them is missed, delayed, or scattered across records, and physicians spend much of their time piecing it together. Existing standards identify where information came from, but they point at conflicts rather than resolve them. We propose TRIAD, a data model that connects the information recorded by nurses and physicians so it can be shared and updated as it changes, supporting the path from a patient's first report to what the physician needs, and allowing an agent to retrieve that information directly rather than search for it. The codebase and artifacts are available under the CC BY-NC-SA 4.0 license at \url{https://github.com/MohamedKhalidmk/TRIAD}.
\end{abstract}

\keywords{Knowledge Graph, RDF, FHIR, Agent Memory, Healthcare Data Privacy}
\maketitle

\section{Introduction}
% Khaled's section: his own text, AI-edited for clarity only (Level C).
% Originals in Appendix A.
When a patient calls a clinic, the information they give may be incomplete or inaccurate, and it is then summarized by a nurse. Errors at this stage have serious consequences. Patients may be prioritized incorrectly, so urgent cases wait; reviews of emergency triage report nurse accuracy between 59.3\% and 82\%~\cite{triage}. Physicians also lose time working from wrong or outdated information in systems that do not present it quickly, spending 49\% of their office day on electronic records and desk work compared with 27\% with patients~\cite{sinsky2016}.

Information about a patient comes from several people and systems over time: the patient's own report of symptoms, the provincial pharmacy dispensing record (PharmaNet)~\cite{pharmanet}, existing clinical records such as imaging and test results, the nurse's structured assessment and recommendation (SBAR)~\cite{sbar}, and the physician's decision. The nurse's recommendation informs the physician's decision; the physician may accept it, reject it, or build on it.

Because this history is created by different people and organizations, records can disagree, and the same patient may appear under different names. When records disagree, the physician cannot easily tell whether a record still reflects what the patient is actually taking~\cite{ismp}, who recorded it and when, who verified it, and whether the nurse's recommendation was accepted, rejected, or built on.

Existing standards such as HL7 FHIR can record where each piece of information came from~\cite{provenance}, but they point at conflicts rather than resolve them. A patient's medication information is scattered across orders, dispenses, administrations, patient reports, and claims~\cite{fhirmeds}, and reconciling it is left to clinicians as a manual process~\cite{smp} that is burdensome and takes time away from direct patient care~\cite{bastola2025}.

\section{Motivation}
Hospitals are slow to use agents in clinical work, for four reasons. First, patient data is in many separate systems. Even where hospitals adopt interoperability standards such as FHIR \cite{fhir}, some data is on paper, and other data is behind an interface that nobody outside the building can reach.

Second, an agent cannot keep all of this data in its conversation window. When the context is long, language models miss relevant facts. The agent does not see that a fact changed, and it uses an old fact.

Third, patient data is sensitive. Personally Identifiable Information (PII) is information that can identify a person, such as a name or a health card number. When PII is linked to health data, it becomes Protected Health Information (PHI). An agent must not show PHI to the wrong person.

Fourth, the use of PHI and PII is subject to strict laws, such as HIPAA in the United States, and PIPEDA and PHIPA in Canada. A hospital cannot use a system that does not obey these laws.

An Electronic Health Record (EHR) is designed for people to read. It does not give the agent a small and explicit view of who asserted what and what replaced what. The agent must infer this at runtime. When the inference is wrong, nobody can find the cause. Wrong answers decrease the trust of staff, and hospitals then do not use agents.

\subsection{Motivating Example}
A patient comes to the emergency room with shortness of breath. At triage, the patient tells the agent that the problem started two days ago. The nurse assesses the patient and records a possible asthma attack. One hour later, the physician examines the patient and revises the assessment to possible heart failure. The physician also tells the agent that the father of the patient died of a heart attack. This fact is not in the EHR yet.

In the evening, the patient asks the agent about the plan for their care. The conversation is now long. The asthma assessment appears many times, and the revision appears once. The agent answers from the old assessment and tells the patient that an inhaler is the plan.

With TRIAD, the agent starts at the patient and finds the current assessment one or two hops away in an RDF \cite{rdf} knowledge graph. It sees that the physician replaced the assessment of the nurse, and it finds the family history. If its answer is wrong, we can trace the error to a specific relationship. In the session of the patient, the agent shows only the data that is approved for the patient. Staff who do not care for this patient get no access to their PHI.

\section{Team Justification}
This project combines the backgrounds of the two team members. Khaled Abuela does his M.Sc. thesis research on memory for language model agents under Professor Shengyao Lu. Asish George has four years of experience with the development of healthcare systems. We selected a problem that we both know from real work, not an artificial one.

\bibliographystyle{ACM-Reference-Format}
\bibliography{bibliography.bib}

% ---------------------------------------------------------------
% APPENDIX: original text before AI-assisted editing (Level C).
% Does not count toward the one-page limit.
% ---------------------------------------------------------------
\clearpage
\appendix
\section{Original Work Before AI-Assisted Editing}
Per the Level C guideline, this appendix contains Mohamed Khaled Abuela's original text for the Abstract and Introduction, written without AI, before AI-assisted editing. AI (Claude) was also used to locate and explain sources.

\subsection{Abstract}
patient wait and goes critical and miss information keeps delaying while physicions is working 24/7 and we are just trying to make standards that arent helping but just pointing triad automates the entire processes from summary to giving the doctor waht he needs using a smart data model that doesnt require the agent to ask questions and fetch but just to pull without reason sharing across nurve and doctor info can be shared and exchanged and editted right awya without any problems all connected

\subsection{Problem Definition Notes}
\textit{Answers to the six problem-definition questions, as originally written:}

1. patient on phone then pharmacy then xrays then ISMP then SBAR

2. patient history and medications different other clinics or out of data, also name different in entities like m.khaled or mohamed khaled or different descriptions of symptoms

3. what patient who wrote it who verified

4. Next patient, conditon, how argent and what they are currently taking and nurse recommendation

5. same

6. why it was recorded who recorded it when and patient full history and condition

\textit{Follow-up answers:} medications the patient currently taking and it gives the doctor an idea and the doctor can use it or reject it or expand it; history and results; human symbols and writing and also what patient and and time and author revision verification etc.

\textit{Problem paragraph:} Patient history can differ as its created by different people and organisation and its not easy to know without wasting the doctors time also a nurse can miss important info over the phone and the doctor can waste sm time on something he didint need to do while also other more important patient suffer

\subsection{Introduction: Opening}
so the patient when he calls a hospital the info he give mgiht not be that accurate or summaries so the nursey summaries it [...] that has some serious problems which are first the ordering of patients are incorrect, causing important patients to wait critical time and sec that the doctor time is wasted with completely wrong info, also doctors deal with the slowest out dated systems that doesnt just give the info right awya we ened a source for that

\subsection{Introduction: Disagreement Between Records}
Because this history is created by different people and organizations, records can disagree, and the same patient may appear under different names. When records disagree, the physician cannot easily tell whether a record still reflects what the patient is actually taking, who recorded it and when, who verified it, and whether the nurse's recommendation was accepted, rejected, or built on.

\subsection{Introduction: Gap}
fhir just signaties what comes from a different source not resolve the conflict its pointing at the problem not offering solvtions and also the sattering thing i saiid and also most of clinical work is still maual and u can find many evidence

\end{document}
